% Owner: classification_e. Root and boundary_d full-source referee PASS.
\begin{proposition}[E.2.14 (Proved): extension of the circle field and exclusion of finite-root collapse]
In the setting of E.2.10, the normalized vector field
$\delta$ of the punctured family extends to a holomorphic
relative vector field on $\mathcal X$ after shrinking
the parameter disc.  Its time-one flow is the identity,
and its real circle action on the central fibre is
equivariantly diffeomorphic to the prescribed circle
action on the punctured fibres.  In particular its
central fixed locus is a smooth rational curve with
normal bundle $\mathcal O(-1)\oplus\mathcal O(-1)$
and the same two opposite circle weights.

The rank-one limit of the normalizing projective
transformations $A_s$ cannot have constant image
$0$ or $1$.  More generally, a collapsed constant
image is impossible whenever the limiting torus's
complex $\delta$ line has two real-linearly
independent periods.  Thus the finite-root
isotrivial possibilities in E.2.12 are excluded.
\end{proposition}
\begin{proof}
Let $\pi:\mathcal X\to\Delta$ be the original
proper holomorphic submersion.  The coherent sheaf
\[
 \mathcal H=\pi_*T_{\mathcal X/\Delta}
\]
is torsion-free: multiplication by a nonzero base
function is injective on the vector bundle upstairs
and hence on its sections.  Its generic rank is
one by E.2.1.  On the smooth disc it is therefore
locally free of rank one.  Choose a generator $V$
after shrinking the disc.  On $\Delta^*$ it has
the form
\[
 V|_{X_s}=b(s)\delta_s,
\]
with $b$ holomorphic and nowhere zero there.
The normalized fields $\delta_s$ agree on overlaps
of local classifying maps: at fixed discrete data
the comparison maps of E.1.3 have fibre-linear
part the identity and fix the marked period
$\widehat\delta$.  Their translations also
preserve the invariant vector field.
Moreover $V$ is vertical for $g$, because this
holds on the dense punctured family.

Choose a regular central target value and a small
bidisc around it on which $g$ is smooth.  The
ordinary or finite-root comparison used in
E.2.12 supplies its smooth marked Jacobian, with
the marked period $\widehat\delta$ extending as
a nonzero holomorphic Lie section.  The resulting
invariant field $\delta_B$ on the actual torsor
is nonvanishing there.  Comparing one nonzero
coefficient of $\delta_B$ with $V$ shows that
$b(s)$ extends holomorphically to $s=0$.
Its independence of the target coordinate follows
from that equality on $s\ne0$ and holomorphic
uniqueness.

If $b(0)=0$, the restriction $V|_{X_0}$ vanishes
on a nonempty open subset of connected $X_0$,
and therefore vanishes everywhere.  In local
coordinates for the submersion $\pi$, every
coefficient of $V$ is then divisible by $s$.
The quotients glue uniquely to a global
holomorphic relative vector field $V/s$.
This contradicts the fact that $V$ generates
the free module $\mathcal H$.  Hence $b(0)\ne0$,
and $V/b(s)$ gives the required global extension
of $\delta$.

The extended field is complete on each compact
fibre of $\pi$.  Its time-one flow is the
identity on the punctured family and hence on
the whole family by holomorphic uniqueness.
Its real flow consequently defines an action
of $\mathbb R/\mathbb Z$ on $\mathcal X$
over the parameter disc.  Average a smooth
metric and an Ehresmann horizontal distribution
over this compact group.  Parallel transport
then gives equivariant diffeomorphisms of
the fibres.  The central action is therefore
the same effective semifree action, with a
single fixed two-sphere, as on the prescribed
fibres.

The fixed set of a holomorphic compact-group
action is a complex submanifold, as follows
from the averaging coordinates used in
E.2.11.  Here it is a proper smooth family
of curves of genus zero, and hence each fibre
is $\mathbb P^1$.  Its two holomorphic normal
weight lines have the same degrees as in
the punctured family.  Those degrees are
$-1,-1$ by C.3.11, or directly by the two
fixed toric edge normals.  The characters
distinguish the lines, so their central
restrictions are exactly
$\mathcal O(-1)\oplus\mathcal O(-1)$.

Suppose now that $A_s$ has rank-one limit
with constant image a finite marked value.
E.2.12 identifies every central regular
torus, with its marked $\delta$ direction,
with the corresponding finite-root torus
at parameter $c(0)$.  The invariant lattice
of $A_j$ has rank two and spans its
one-dimensional complex fixed subspace.
It contains $\widehat\delta$; its second
primitive period is independent over
$\mathbb R$.  The complex fixed subspace
is precisely the $\delta$ line, since
the other complex eigenvalue is a
primitive third or fourth root of unity.
Consequently, in the normalization in
which the $\delta$ period is $1$, there
is a constant second period
\[
 \kappa\in\mathbb C\setminus\mathbb R.
\]
This is a period of the invariant field
on every central regular fibre.  The
time-$\kappa$ flow is therefore the
identity on their dense union, and
hence on all of $X_0$.

The flow of the nonzero field $\delta_0$
thus factors through a nontrivial
holomorphic action of the compact
elliptic curve
$E=\mathbb C/(\mathbb Z+\kappa\mathbb Z)$.
Such an action has no fixed point on
a connected complex manifold unless
it is trivial.  Indeed its derivative
at a fixed point is a holomorphic
representation of the compact complex
group into a general linear group,
so is constant.  Averaging local
coordinates over the compact group
then makes the action trivial near
that point, and holomorphic uniqueness
makes it trivial everywhere.

The elliptic action here is nontrivial
because $\delta_0\ne0$.  It therefore
has no fixed point.  This contradicts
the fixed rational curve of the
extended real circle action: a zero
of the generating field is fixed
by its entire complex flow.  Equivalently,
the resulting nowhere-zero real vector
field would give $\chi(X_0)=0$,
whereas $X_0$ is diffeomorphic to $S^6$.
The same argument applies to an ordinary
constant-image torus if its $\delta$
line has a second independent period.
\end{proof}

\begin{remark}[What the elliptic-action argument does not assume]
The argument uses the global one-dimensional
space of vector fields on the known punctured
members and the actual marked period comparison
on a central regular bidisc.  It does not
assume that all translations of an ordinary
central torus extend over the singular fibre.
For an ordinary torus the $\delta$ line may
have only the one specified real period;
such a case is not excluded by the
elliptic-action argument alone.
\end{remark}
